\begin{center}
    \textbf{ВВЕДЕНИЕ}
\end{center}
\addcontentsline{toc}{chapter}{ВВЕДЕНИЕ}

Возможность гибко задавать часто используемые действия системы напрямую влияет на удобство взаимодействия пользователя с компьютером. Современные геймпады, такие как Xbox Series S, предоставляют широкие возможности для альтернативного способа управления компьютером, включая функции мыши. Это особенно актуально в ситуациях, когда использование традиционной мыши затруднено или невозможно, например при подключении компьютера к телевизору, в домашних кинотеатрах или при работе с медиацентрами.

Геймпад Xbox Series S, благодаря своим точным аналоговым стикам и эргономичной конструкции, идеально подходит для реализации функций мыши. Разработка драйвера, позволяющего использовать данный контроллер в качестве полноценного устройства ввода типа "мышь", является актуальной задачей, которая расширяет функциональность контроллера и повышает удобство работы пользователя в различных сценариях.

Данная работа посвящена разработке драйвера для геймпада Xbox Series S, который преобразует контроллер в устройство мыши с поддержкой перемещения курсора, нажатия кнопок и прокрутки.